\documentclass[11pt,a4paper,roman]{moderncv}
\moderncvstyle{banking}
\moderncvcolor{purple}
\nopagenumbers{}
\usepackage[utf8]{inputenc}
\usepackage[T1]{fontenc}
\usepackage{lmodern}
\usepackage[scale=0.96,top=0.48cm,bottom=0.48cm,left=0.72cm,right=0.72cm]{geometry}
\renewcommand{\baselinestretch}{0.92}
\setlength{\hintscolumnwidth}{2.20cm}
\setlength{\footskip}{18pt}

\name{Wahid}{SAMER}
\address{Cergy (95), France}{}{}
\mobile{07 53 10 95 74}
\email{samer.ahmed.wahid@gmail.com}
\social[linkedin]{Wahid SAMER}
\social[github]{wahidrt}
\title{Alternance M2 -- Weather \& Energy Trader / Quantitative Trading}

\begin{document}
\makecvtitle
\vspace{-3.4em}
\begin{center}
\parbox{0.97\textwidth}{\centering\small\itshape
Étudiant en \textbf{Master Mathématiques Appliquées à l'Ingénierie Financière}, orienté \textbf{trading quantitatif, modélisation stochastique et gestion des risques}. Expérience académique en \textbf{Python/VBA, séries temporelles, pricing, backtesting et modèles quantitatifs}. Recherche d'une \textbf{alternance M2 de 12 mois à partir du 21 septembre 2026}.}
\end{center}
\vspace{0.15em}

\section{Formation}
\cventry{2025 -- 2027}{Master Mathématiques Appliquées à l'Ingénierie Financière}{CY Cergy Paris Université / CY Tech}{Cergy}{}{\textbf{M1 :} processus stochastiques, probabilités, simulation stochastique, EDP, méthodes numériques avancées, optimisation, gestion de portefeuille, Python/VBA.\\
\textbf{M2 2026--2027 :} stochastic calculus, séries temporelles, machine learning avec Python, model calibration and simulation, portfolio management, fixed income.}
\cventry{2021 -- 2025}{Licence Mathématiques et Applications}{Université d'Évry Paris-Saclay}{Évry}{}{Probabilités, mesure et intégration, analyse, algèbre, topologie, analyse numérique et programmation.}

\section{Projets ciblés Trading, Risque \& Pricing}
\cventry{2026}{Python / Machine Learning}{Robot de trading quantitatif adaptatif}{Projet M1 -- équipe de 4}{}{Modélisation de séries financières multi-actifs : cointégration, VECM, filtre de Kalman, processus d'Ornstein--Uhlenbeck et détection de régimes ; développement d'une architecture de backtesting et d'analyse quantitative.}
\cventry{2026}{Python}{Backtesting, portefeuille et risque}{Projet M1}{}{Conception de modules de portefeuille et P\&L, exécution simulée à $t+1$, coûts de transaction, equity curve, Sharpe, Sortino, drawdown et turnover.}
\cventry{2026}{Python / \LaTeX}{Pricing d'options -- Black--Scholes \& Heston}{Mémoire M1}{}{Calcul stochastique, Itô, EDP, Feynman--Kac, pricing Call/Put, Greeks, méthodes numériques et étude de la volatilité stochastique de Heston.}
\cventry{2025}{Excel / VBA}{Pricer d'options vanilles et stress tests}{Projet M1}{}{Valorisation Black--Scholes, calcul de $\Delta,\Gamma,\Theta,\nu,\rho$ et scénarios de stress sur la volatilité.}

\section{Compétences}
\cvitem{Quant / Trading}{Séries temporelles, cointégration, VECM, Kalman, Ornstein--Uhlenbeck, backtesting, gestion de portefeuille, mesures de risque.}
\cvitem{Pricing / Risque}{Black--Scholes, Heston, Greeks, simulation stochastique, EDP, Feynman--Kac, stress tests.}
\cvitem{Programmation}{Python (NumPy, Pandas, SciPy, Scikit-learn, Matplotlib), VBA, C/C++, SQL ; Excel, Git/GitHub, \LaTeX.}
\cvitem{Langues}{Français : courant ; Arabe : langue maternelle ; Anglais : intermédiaire.}

\section{Expérience \& Engagement}
\cventry{Mai -- juin 2025}{Professeur de mathématiques -- stagiaire}{Lycée Pierre Mendès France}{Vitrolles}{}{Préparation et transmission de notions de Terminale ; rigueur, pédagogie et communication scientifique.}
\cventry{2019}{Technicien aéronautique -- stagiaire}{G1 Aviation}{Gap}{}{Assemblage/maintenance dans un environnement exigeant en précision, procédures et sécurité.}
\cvitem{Associatif}{Bénévole aux \textbf{Restaurants du Cœur} pendant près de 3 ans : solidarité, accueil et travail d'équipe.}

\end{document}
